\section{Introduction}
Statistical watermarks over text are being deployed to attribute leaked
documents and flag AI-generated content. Vendors and papers quote detection
rates of 90--100\% under calibration, and deployments inherit these numbers
as operating expectations. This paper shows that three systematic failure
modes---each easy to trigger with ordinary text, each invisible to the
standard evaluation protocol (same-domain calibration, random test
splits)---collapse those numbers on real corpora, and that two of the three
have no fix inside the watermark itself.

We study \emph{lexical statistical watermarks}: the verifier applies a
keyed, deterministic predicate to words, accumulates per-band vote counts,
and thresholds a $z$-style existence score, optionally decoding a multi-user
identifier from band signs. Our host system, AAWM (v0.15.0), is one instance
of this class; the pitfalls, however, are properties of the statistic, not
of the implementation. Where a claim could be implementation-specific we
reproduce it on (i) an independent 80-line reference implementation with no
shared code, and (ii) the official implementation of PostMark
\cite{postmark}, a post-hoc black-box watermark of a different statistical
family.

\paragraph{Finding 1: correlated votes.}
The word predicate is a deterministic function of the key. Repeating one
dictionary word $N$ times therefore contributes $N$ perfectly correlated
``samples'' to a test whose binomial variance formula assumes independence.
The existence score grows $\propto\sqrt{N}$ without bound: an unwatermarked
form letter that repeats a single in-dictionary word 640 times scores
$25.3$ against a default threshold of $8$ (Figure~\ref{fig:eb1}). Capping
votes per word per document at 3 bounds the attack ($\le 1.73$) and, on
in-domain text, is also the \emph{best operating point}: highest true-positive
rate with zero false positives (Figure~\ref{fig:eb2}). The reference
implementation reproduces the identical curve (Figure~\ref{fig:eb4}).

\paragraph{Finding 2: non-transferable null models.}
Thresholds fitted on a calibration corpus are domain-specific in a strong,
measured sense: calibrating on technical prose and testing on classical
Chinese novels or English literature raises the false-positive rate from
0--6.9\% (uncalibrated) to 3.5--84\%---\emph{worse than not calibrating at
all} (Table~\ref{tab:eb3}). Only in-domain calibration restores FPR to
0--7.7\%, and even that requires an adequate corpus size: with only 14
calibration documents, in-domain calibration itself fails (46.7\% FPR).
PostMark's bounded presence-ratio statistic is measurably more
domain-stable (0.5--5\% cross-domain FPR), showing that pitfall severity is
statistic-dependent.

\paragraph{Finding 3: the template attack.}
Texts built from repeated template sentences (form letters, mail merges,
review-boosting content) produce high existence scores \emph{and} false
attribution to specific registered users---13--28\% across two independent
samples---because the noise genuinely aligns with some candidate's bit
pattern (candidate score ratio up to 0.64, band-alignment rate up to 0.81;
Figure~\ref{fig:eb6}). Abstain-on-low-margin protocols do not prevent this:
the alignment is real, so the margin is large. On natural templates with
slot filling, false attribution disappears (0\%) but existence false
positives remain high without calibration (30\%); only in-domain calibration
or a business-layer similarity prefilter mitigates (Table~\ref{tab:eb7}).
PostMark's presence ratio is structurally immune to this attack (score
exactly 0 on 200 template texts).

\paragraph{Contributions.}
\begin{itemize}
\item \textbf{C1} We construct, reproduce on two independent implementations,
and bound the correlated-vote inflation, and characterize the cap trade-off
(\S\ref{sec:f1}).
\item \textbf{C2} We provide the first cross-domain measurement of empirical
threshold transfer: cross-domain calibration is a net negative; corpus size
matters as much as domain match (\S\ref{sec:f2}).
\item \textbf{C3} We identify the template attack as a structural boundary
of statistical attribution, characterize its mechanism, and rank mitigations
(\S\ref{sec:f3}).
\item \textbf{C4} We show the pitfall profile is statistic-dependent
(presence-ratio vs.\ $z$-test), turning our findings into a deployment
checklist (\S\ref{sec:recs}).
\end{itemize}

All experiments are fixed-seed and reproducible from the released artifact
(AAWM v0.15.0, experiment scripts, and data files).
